\clearpage\section*{TÓM TẮT}\addcontentsline{toc}{section}{Tóm tắt}
Nghiên cứu đề xuất mô hình tối ưu lịch sơ tán cho tòa nhà A4, Đại học Phenikaa, trên mạng có hướng với ràng buộc thời gian, khả năng thông hành và nguy cơ động. Miền tính toán gồm bảy tầng, 27 khu phòng và tổng sức chứa trong tập dữ liệu là 1.360 người. Hai cụm cầu thang và các sảnh tiếp cận được biểu diễn tường minh; khu vệ sinh nằm ngoài miền mô phỏng. Hành lang tầng 1 được xét là vùng hở nối trực tiếp với không gian ngoài. Vì vậy, cửa phòng tầng 1 và vị trí tiếp cận hành lang từ sảnh thang được mô hình hóa thành các nút thoát hấp thụ.

Phương pháp M4 là bộ điều phối luồng toàn cục trên mạng mở rộng theo thời gian (TEN), kết hợp Dijkstra đa nguồn, Dinic và luồng chi phí nhỏ nhất. Ba phương pháp M1--M3 lần lượt sử dụng tuyến cố định, tái định tuyến theo nguy cơ và lựa chọn tham lam theo năng lực. Hàm mục tiêu của M4 được xác lập theo thứ tự ưu tiên: tối đa số người thoát trong thời hạn; tối thiểu thời điểm thoát cuối của lượng người này; và tối thiểu tổng thời điểm thoát. Luồng nguyên được phân rã thành lịch di chuyển và chờ, sau đó kiểm tra độc lập về tính liên tục, sức chứa, năng lực cạnh, điều kiện an toàn và bảo toàn dân số.

Thực nghiệm gồm 16 tổ hợp cơ sở, đối chứng chia đều hai thang, khảo sát hệ số M3, năm lượt đo lặp M4, 15 phân bố dân số và các khảo sát độ nhạy. M4 sơ tán đủ 1.360 người trong 390 s khi không có sự cố, so với 500 s của tuyến cố định và 395 s của đối chứng chia đều hai thang. Cận dưới theo lát cắt cũng bằng 390 s. Khi khói tại phòng 403, M4 hoàn tất trong 425 s; khi khói lan từ cầu thang, lượng thoát tối đa trong 900 s là 1.248, tương ứng 91,76\%.

Kết quả cho thấy điều phối tuyến và thời điểm khởi hành có tác động đồng thời đến lượng người thoát và mức phân bố tải giữa hai thang. Tính tối ưu được giới hạn trong mô hình rời rạc với lịch nguy cơ biết trước. Các tham số thông hành và lan truyền nguy cơ chưa được hiệu chỉnh bằng số đo hiện trường; kết quả phục vụ đánh giá thuật toán và cần được kiểm chứng thêm trước khi ứng dụng vào vận hành công trình.

\noindent\textbf{Từ khóa:} sơ tán; đồ thị động; mạng thời gian; luồng cực đại; luồng chi phí nhỏ nhất; kiểm chứng lịch.
